\section{The 2-adic weights}\label{sec:fibers}
We now construct the measurable family of shifts. Their logarithmic weights must satisfy the criterion of Section~\ref{sec:criterion} while changing in opposite directions at the two ends when the base moves by one.

\subsection{The base transformation}
For the inverse-limit base $X=\Z_2$ fixed in the introduction and $m\ge1$, let
\[
 N_m=2^m,\qquad T_m(x)=x\bmod N_m\in\{0,\ldots,N_m-1\}.
\]
Thus each residue is uniformly distributed. The map $\theta x=x+1$ is an invertible probability-preserving transformation, and
\begin{align}
 T_{m+1}(x)-T_m(x)&\in\{0,N_m\},\label{eq:residue-nesting}\\
 T_m(\theta x)-T_m(x)&=
 \begin{cases}1,&T_m(x)<N_m-1,\\-(N_m-1),&T_m(x)=N_m-1.\end{cases}\label{eq:residue-step}
\end{align}
The finite residue classes generate the Borel sigma-algebra of $X$.

\begin{lemma}\label{lem:ergodicity}
The transformation $\theta$ is ergodic.
\end{lemma}
\begin{proof}
If $E$ is invariant modulo null sets, $\theta$ cyclically permutes the $N_m$ residue classes, so
\[
 \mu(E\cap\{T_m=r\})=\frac{\mu(E)}{N_m}\qquad(0\le r<N_m).
\]
Thus $E$ is independent of every finite union of residue classes. By the Monotone Class Theorem it is independent of every measurable set, in particular itself. Hence $\mu(E)=\mu(E)^2$.
\end{proof}

\subsection{The weights}
Fix
\[
 c=10,\qquad D=100,\qquad \ell_0=1,\quad\ell_m=8^m\ (m\ge1).
\]
The integer blocks $I_m=\{k:\ell_{m-1}\le k<\ell_m\}$ partition the positive integers. Set $J_1=0$ and $J_{m+1}=J_m+DN_m$; equivalently $J_m=D(N_m-2)$. Define $a_0^+(x)=a_0^-(x)=c$, and for $k\in I_m$ put
\begin{align}
 a_k^+(x)&=c(k+1)+J_m+DT_m(x),\label{eq:a-plus}\\
 a_k^-(x)&=c(k+1)+J_m+D(N_m-1-T_m(x)).\label{eq:a-minus}
\end{align}
On the full integer axis define
\begin{equation}\label{eq:bilateral-a}
 a_j(x)=\begin{cases}a_j^+(x),&j\ge0,\\a_{-j-1}^-(x),&j<0,\end{cases}
 \qquad \beta_j(x)=e^{-a_j(x)}.
\end{equation}
The two central weights satisfy $\beta_{-1}(x)=\beta_0(x)=e^{-c}$ for every $x$.
The linear terms in $a_k^\pm$ give quadratic growth of the cumulative logarithms. The smaller residue corrections change by $+D$ and $-D$ under $\theta$ unless the residue wraps around; the increments of $J_m$ preserve monotonicity across block boundaries. We first check the invariant-subspace criterion, then use those opposite changes to construct the backward intertwiner.

\begin{proposition}\label{prop:fibers}
For every $x\in X$, the operator $S_xe_j=\beta_j(x)e_{j+1}$ is bounded, has norm $e^{-c}$, and has precisely the standard tails as its closed invariant subspaces. For each $v\in K$, the map $x\mapsto S_xv$ is measurable.
\end{proposition}
\begin{proof}
We verify positivity, regularity, and growth separately. The corrections to $c(k+1)$ in \eqref{eq:a-plus}--\eqref{eq:a-minus} are nonnegative. Thus
\begin{equation}\label{eq:beta-sum}
 0<\beta_j(x)\le e^{-c},\qquad
 \sum_{j\in\Z}\beta_j(x)\le2\sum_{k\ge0}e^{-c(k+1)}<\infty.
\end{equation}
Within a block, both $a_k^\pm$ increase by $c$ at each step. At a boundary from $I_m$ to $I_{m+1}$, put $\delta=T_{m+1}(x)-T_m(x)\in\{0,N_m\}$. The increments are, respectively,
\[
 c+DN_m+D\delta\ge c,
 \qquad c+D(N_{m+1}-\delta)\ge c.
\]
Also $a_1^\pm\ge2c>a_0^\pm$. Consequently $\beta_j$ is nonincreasing for $j\ge0$ and nondecreasing for $j\le0$, including the equality at $-1,0$.

Set $w_0(x)=1$ and define $w_{j+1}(x)/w_j(x)=\beta_j(x)$ for all $j\in\Z$. This gives a positive decreasing sequence. Since $\log w_{j+1}-\log w_j=\log\beta_j$, its logarithm has the required convexity and concavity.

For $k\in I_m$ we have $N_m\le2k^{1/3}$, and hence
\begin{equation}\label{eq:a-error}
 0\le a_k^\pm(x)-c(k+1)\le4Dk^{1/3}\qquad(k\ge1).
\end{equation}
The bound is uniform in $x$. Summing gives, as $n\to\infty$,
\begin{align}
 \log w_n(x)&=-\sum_{k=0}^{n-1}a_k^+(x)
   =-\frac{cn(n+1)}2+O(n^{4/3}),\label{eq:w-positive}\\
 \log w_{-n}(x)&=\sum_{k=0}^{n-1}a_k^-(x)
   =\frac{cn(n+1)}2+O(n^{4/3}).\label{eq:w-negative}
\end{align}
Thus \eqref{eq:quadratic} holds with $a=c/2=5>\log3$. Moreover the lower bounds $a_k^\pm\ge c(k+1)$ imply
\[
 w_n(x)^{1/n}\le e^{-c(n+1)/2},\qquad
 w_{-n}(x)^{-1/n}\le e^{-c(n+1)/2}\qquad(n\ge1).
\]
These bounds prove \eqref{eq:root-sum}. Lemma~\ref{lem:domar} therefore applies to every fiber.

The norm statement follows from \eqref{eq:beta-sum} and $\beta_0=e^{-c}$. Each coefficient $\beta_j$ is measurable, so $x\mapsto S_xe_j$ is measurable. Finite basis expansions and the uniform norm bound give strong measurability on every vector of $K$.
\end{proof}

We may now define $S\in\B(H_0)$ by $(Sf)(x)=S_xf(x)$. It has norm $e^{-c}$. For each $g\in L^\infty(X)$, the scalar multiplier $(M_gf)(x)=g(x)f(x)$ commutes with $S$. Before constructing the operator that links different fibers and moves their tails backward, we record the uniform power bound for $S$.

\begin{proposition}\label{prop:normquasi}
The operator $S$ is nonzero and norm-quasinilpotent. More precisely, for
all integers $n\ge1$,
\[
 \|S^n\|\le \exp\left(-c\left\lfloor\frac{(n+1)^2}{4}\right\rfloor\right).
\]
\end{proposition}
\begin{proof}
Define $d(j)=j+1$ for $j\ge0$ and $d(j)=-j$ for $j<0$. The weight definitions
imply $a_j(x)\ge c d(j)$ for every $x,j$. Since
\[
 S_x^n e_j=\exp\left(-\sum_{r=0}^{n-1}a_{j+r}(x)\right)e_{j+n},
\]
the norm of each fiber power is bounded by the claimed quantity. Indeed the
minimum of $\sum_{r=0}^{n-1}d(j+r)$ over $j\in\mathbb Z$ is attained by an
interval centered as closely as possible on the two minima at $-1$ and $0$.
For $n=2m$ this minimum is $m(m+1)$; for $n=2m+1$ it is $(m+1)^2$.
Both values are $\lfloor(n+1)^2/4\rfloor$. The bound is independent of $x$,
so integrating the fiber estimate proves it for $S^n$ itself. Consequently
$\lim_{n\to\infty}\|S^n\|^{1/n}=0$. Finally
$S(1\otimes e_0)=e^{-c}(1\otimes e_1)\ne0$.
\end{proof}
